% =============================================================================
% Title:        Flow Controller DB15 to DB9 Connector Diagram
% Author:       Nathaniel Thomas
% Affiliation:  Texas A&M University, Department of Nuclear Engineering
% Email:        nathaniel@tamu.edu
% Date Created: 3/7/2026
% Version:      1.0
%
% Description:
%   A colored wiring diagram of the FLiNaK purifier flow controllers.
%   This diagram shows the wire connections for the DB15 female to DB9 female
%   connecting cables used to connect the Aalborg flow controllers to the
%   NI USB-6353 that controls the purifier.
%
% =============================================================================

\begin{circuitikz}
    \node[draw, minimum width=1.5cm, minimum height=4.5cm] (A) {DB9};
    % Right DB9 connector
    \node[draw, minimum width=1.5cm, minimum height=7.5cm, right=6cm
    of A] (B) {DB15};

    % Left connector pins
    \foreach \i [evaluate=\i as \y using {2 - 0.5*(\i-1)}] in {1,...,9}
    {
      \node[circle, draw, inner sep=1pt] (L\i) at (1.5,\y) {};
      \node[left] at (L\i) {\i};
    }
    \node[circle, draw, inner sep=1pt] (LGND) at (0,-3) {};
    \node[above] at (LGND) {GND};

    % Right connector pins
    \foreach \i [evaluate=\i as \y using {3.5 - 0.5*(\i-1)}] in {1,...,15}
    {
      \node[circle, draw, inner sep=1pt] (R\i) at (6,\y) {};
      \node[right] at (R\i) {\i};
    }
    \node[circle, draw, inner sep=1pt] (RGND) at (7.5,-4.5) {};
    \node[above] at (RGND) {GND};

    \newcommand{\pinconnect}[3]{%
      % #1 = DB9 pin, #2 = DB15 pin, #3 = color
      \draw[#3, thick] (L#1) -- (R#2);
    }

    \pinconnect{1}{1}{orange}
    \pinconnect{2}{2}{yellow}
    \pinconnect{3}{4}{green}
    \pinconnect{4}{8}{blue}
    \pinconnect{5}{7}{red}
    \pinconnect{6}{10}{violet}
    \pinconnect{7}{11}{darkgray}
    \pinconnect{8}{12}{lightgray}
    \pinconnect{9}{5}{black}
    \draw[brown, thick] (LGND) -| (L9);
    \draw[brown, thick] (RGND) -| (R15);
    \draw[brown, thick] (R15.west) -- ++(-1,0) |- (R5.west);
    \draw[brown, thick] (R3.west) -- ++(-1,0) |- (R5.west);
\end{circuitikz}
